\documentclass[10pt,twocolumn]{article}
\usepackage[margin=0.8in]{geometry}
\usepackage{times}
\usepackage{amsmath}
\usepackage{booktabs}
\usepackage{natbib}

\title{Confidence-Gated Early Exit for On-Device Speech Recognition}
\author{Anonymous Authors}
\date{}

\begin{document}

\maketitle

\begin{abstract}
Speech recognisers on phones and watches must answer within tens of milliseconds while sharing a small battery with everything else.
We add exit heads to a Conformer encoder and stop computing as soon as a CTC head is confident about the current chunk of audio.
On LibriSpeech, the gated model keeps the word error rate of the full model within 0.2 points on test-clean while running 38\% fewer encoder layers on average.
\end{abstract}

\section{Introduction}

On-device speech recognition keeps audio private and works without a network, but the device sets a hard budget for latency and energy.
Most utterances are easy: clear speech, common words and little noise.
A recogniser that spends the same computation on every chunk of audio therefore wastes most of it.

Early exit \citep{teerapittayanon2016branchynet} attaches small classifiers to intermediate layers and stops when one of them is confident.
It is widely used for image classification but less so for streaming speech recognition, where the output is a sequence and the decision must be made chunk by chunk.
We show that a simple confidence gate on CTC heads \citep{graves2006connectionist} works well for a Conformer encoder \citep{gulati2020conformer}.

\section{Related Work}

Early exit was first studied for image classifiers, where each exit head predicts the class and the network stops at the first confident head \citep{teerapittayanon2016branchynet}.
For sequence outputs the decision is harder, since one confident frame says little about the rest of the utterance.
Streaming recognisers process audio in chunks, which gives a natural unit for the decision: the whole chunk exits at the same layer.
Our gate follows this design and uses the CTC posteriors \citep{graves2006connectionist} that the heads already compute, so it adds no parameters.

\section{Method}

The encoder has $L = 16$ Conformer layers.
After layers 4, 8 and 12 we add a linear CTC head that predicts a distribution $q_\ell(\cdot \mid t)$ over tokens for every frame $t$ of the current chunk.
We exit after layer $\ell$ if the mean maximum probability over the chunk exceeds a threshold:
\begin{equation}
\frac{1}{|C|} \sum_{t \in C} \max_{v} \, q_\ell(v \mid t) \;\ge\; \gamma,
\label{eq:gate}
\end{equation}
where $C$ is the set of frames in the chunk.
All heads are trained jointly with the sum of their CTC losses.

\section{Experiments}

We train on the 960 hours of LibriSpeech \citep{panayotov2015librispeech} and report the word error rate (WER) on test-clean and test-other.
Chunks are 640\,ms long, and $\gamma = 0.9$ unless stated otherwise.

\begin{table}[t]
\centering
\caption{WER (\%) on LibriSpeech and mean number of encoder layers run per chunk.}
\label{tab:wer}
\begin{tabular}{lccc}
\toprule
Model & clean & other & Layers \\
\midrule
Full (16 layers) & 3.1 & 7.4 & 16.0 \\
Static exit at 8 & 4.6 & 10.9 & 8.0 \\
Gated, $\gamma = 0.9$ & 3.3 & 7.9 & 9.9 \\
\bottomrule
\end{tabular}
\end{table}

Table~\ref{tab:wer} shows that the gate keeps most of the accuracy of the full model while running far fewer layers.
On test-other the gap is larger, because noisy chunks exit late or not at all.
The gated model is also faster than every other on-device recogniser.

\paragraph{Effect of the threshold.}
Lowering $\gamma$ trades accuracy for speed.
With $\gamma = 0.8$ the model runs 7.6 layers per chunk on average and reaches 3.6\% WER on test-clean; with $\gamma = 0.95$ it runs 12.1 layers and reaches 3.2\%.
Between these values the WER changes smoothly, so the threshold can be set at run time to match the battery level of the device.

\paragraph{Where do chunks exit?}
On test-clean, 41\% of the chunks exit after layer 4, 27\% after layer 8, 12\% after layer 12, and the remaining 20\% run all 16 layers.
Chunks that run the full encoder are mostly silence followed by speech, rare names, and overlapping speakers.
Silence alone exits early, since the blank token dominates every frame.

\paragraph{Latency.}
On a mid-range phone CPU, a 640\,ms chunk takes 31\,ms with the full encoder and 19\,ms on average with the gate, including the cost of the exit heads.

\section{Conclusion}

A confidence gate on intermediate CTC heads makes a Conformer encoder cheaper on easy audio without retraining it for every budget.

Our gate has two limitations.
It decides per chunk, so a chunk that is easy at the start and hard at the end runs to the layer its hard part needs.
And we have tested it only on read English speech; spontaneous and accented speech may exit much later.

\begin{thebibliography}{4}
\bibitem[Graves et~al.(2006)]{graves2006connectionist}
Alex Graves, Santiago Fern{\'a}ndez, Faustino Gomez, and J{\"u}rgen Schmidhuber.
\newblock Connectionist temporal classification: labelling unsegmented sequence data with recurrent neural networks.
\newblock In \emph{Proceedings of the 23rd International Conference on Machine Learning}, 2006.

\bibitem[Gulati et~al.(2020)]{gulati2020conformer}
Anmol Gulati, James Qin, Chung-Cheng Chiu, Niki Parmar, Yu~Zhang, Jiahui Yu, Wei Han, Shibo Wang, Zhengdong Zhang, Yonghui Wu, and Ruoming Pang.
\newblock Conformer: Convolution-augmented transformer for speech recognition.
\newblock In \emph{Proceedings of Interspeech}, 2020.

\bibitem[Panayotov et~al.(2015)]{panayotov2015librispeech}
Vassil Panayotov, Guoguo Chen, Daniel Povey, and Sanjeev Khudanpur.
\newblock Librispeech: an {ASR} corpus based on public domain audio books.
\newblock In \emph{IEEE International Conference on Acoustics, Speech and Signal Processing (ICASSP)}, 2015.

\bibitem[Teerapittayanon et~al.(2016)]{teerapittayanon2016branchynet}
Surat Teerapittayanon, Bradley McDanel, and H.~T. Kung.
\newblock {BranchyNet}: Fast inference via early exiting from deep neural networks.
\newblock In \emph{International Conference on Pattern Recognition (ICPR)}, 2016.
\end{thebibliography}

\end{document}
